% ID-5-02 (TO) · Mathematical Foundations of Diffusion Models · IIT Jammu
% HW 2 — Distributions and Maximum Likelihood
% Build:  latexmk -pdf hw2.tex     (or: pdflatex hw2.tex)
\documentclass[11pt,a4paper]{article}

\usepackage[margin=2.4cm]{geometry}
\usepackage{amsmath,amssymb}
\usepackage{booktabs}
\usepackage[colorlinks=true,linkcolor=black,urlcolor=blue]{hyperref}

\newcommand{\pts}[1]{\hfill\textbf{(#1pts)}}
\newcommand{\pt}[1]{\hfill\textbf{(#1pt)}}
\newcommand{\total}[1]{\par\medskip\noindent\textbf{Total: #1 points}}
\newcommand{\PP}{\mathbb{P}}
\newcommand{\EE}{\mathbb{E}}
\newcommand{\Var}{\operatorname{Var}}
\newcommand{\sig}{\operatorname{\sigma}}

% (a) (b) (c) labelling, without enumitem
\renewcommand{\theenumi}{\alph{enumi}}
\renewcommand{\labelenumi}{(\theenumi)}
\newenvironment{parts}
  {\begin{enumerate}\setlength{\itemsep}{0.45em}\setlength{\parskip}{0pt}}
  {\end{enumerate}}

\begin{document}

\begin{center}
  {\Large\bfseries HW 2 — Distributions and Maximum Likelihood}\\[0.4em]
  {\small Mathematical Foundations of Diffusion Models · Monsoon 2026 · IIT Jammu}
\end{center}

\medskip
\begin{center}
\begin{tabular}{@{}ll@{}}
\toprule
\textbf{Released} & Monday 10 August 2026 \\
\textbf{Due}      & Tuesday 25 August 2026, in class \\
\textbf{Marks}    & 50 points --- graded on correctness \\
\bottomrule
\end{tabular}
\end{center}

\medskip
\noindent
\textbf{Graded on correctness.} Unlike HW 1, this set is marked on the work
itself. Marks are carried by the \emph{reasoning} as much as the final answer: a
correct number with no derivation earns very little, and a wrong number reached
by a sound argument earns most of the credit. Show every step you actually used.
Solutions are posted after the deadline.

\bigskip
\hrule
\medskip

\noindent
\textbf{Notation used throughout.} Three things this sheet leans on.

\begin{itemize}\setlength{\itemsep}{0.4em}
  \item \textbf{i.i.d.} — \emph{independent and identically distributed}. Saying
        $x_1,\dots,x_n$ are i.i.d.\ from some distribution means each was drawn
        from the \emph{same} distribution, and no draw influenced any other. It is
        what lets you multiply the individual probabilities together to get the
        probability of the whole sample.
  \item \textbf{Density vs likelihood} — the same formula, read in two
        directions. Fix the parameters and let the data vary and you have a
        \emph{density} (or pmf): it describes how data comes out, and it
        integrates or sums to $1$. Fix the observed data and let the
        \emph{parameters} vary and you have the \emph{likelihood}: a score saying
        how well each candidate parameter explains what you actually saw. It does
        \emph{not} integrate to $1$ over the parameters, and it is not a
        distribution over them. Maximum likelihood means: pick the parameter with
        the best score.
  \item \textbf{$\Gamma$ and $B$} (Problem 1 only) — $\Gamma$ is the factorial
        extended to non-integers, and the only property you need is
        $\Gamma(z+1) = z\,\Gamma(z)$. The Beta function
        $B(\alpha,\beta) = \Gamma(\alpha)\Gamma(\beta)/\Gamma(\alpha+\beta)$ is
        just the constant that makes the Beta density integrate to $1$.
\end{itemize}

\medskip
\noindent
Show your working. Where you set a derivative to zero, say why the critical point
is a maximum and not a minimum.

\bigskip
\hrule
\bigskip

\section*{Problem 1}

A seed company sells lots of 500 seeds. Not every seed germinates, and the
\emph{germination rate} — the fraction of a lot that sprouts — is not the same
number for every lot: soil, storage and harvest conditions all vary, so across
lots the rate itself scatters somewhere in $[0,1]$.

To model a quantity that is confined to $[0,1]$ and varies, the standard choice is
the \textbf{Beta} distribution. For shape parameters $\alpha, \beta > 0$ its
density is

\[
  f(x; \alpha, \beta) \;=\; \frac{1}{B(\alpha,\beta)}\,
  x^{\alpha - 1}\,(1-x)^{\beta - 1},
  \qquad 0 < x < 1,
\]

and $f = 0$ outside $[0,1]$, where

\[
  B(\alpha,\beta) \;=\; \int_0^1 t^{\alpha-1}(1-t)^{\beta-1}\,dt
  \;=\; \frac{\Gamma(\alpha)\,\Gamma(\beta)}{\Gamma(\alpha+\beta)} .
\]

Larger $\alpha$ pushes the mass towards $1$, larger $\beta$ towards $0$, and both
large together concentrates it. What the company wants to quote is the average
germination rate and how much it varies from lot to lot — so, the mean and the
variance.

Compute both \emph{from the definition of expectation}: integrate against the
density. Do not quote a formula.

\begin{parts}
  \item Show that $f$ integrates to $1$. (One line, given the definition of
        $B$ above.) \pt{1}
  \item Compute $\EE[X] = \int_0^1 x\,f(x;\alpha,\beta)\,dx$. The trick
        throughout: multiplying by $x$ turns the integrand into another Beta
        kernel with $\alpha$ shifted, so the integral is a ratio of Beta
        functions. Then use $\Gamma(z+1) = z\Gamma(z)$. \pts{3}
  \item Compute $\EE[X^2]$ the same way. \pts{3}
  \item Deduce $\Var(X) = \EE[X^2] - (\EE[X])^2$ and simplify to a single
        fraction. \pts{2}
  \item Check your mean at $\alpha = \beta$, and say why that value is obvious
        from the shape of the density without doing any integration. \pt{1}
\end{parts}

\total{10}

\bigskip
\hrule
\bigskip

\section*{Problem 2}

Probabilities are awkward to work with directly: they are trapped in $[0,1]$, so
you cannot add to one without risking leaving the interval. The standard fix is
to move to the \textbf{log-odds} scale. If an event has probability $u$, its odds
are $u/(1-u)$ and its log-odds are

\[
  \operatorname{logit}(u) \;=\; \log\!\frac{u}{1-u},
\]

which stretches $(0,1)$ out onto the whole real line. This is the \emph{inverse}
of the sigmoid $\sig(z) = 1/(1+e^{-z})$, the function that has been squashing
scores into probabilities all through this course.

Here is the question. Suppose a probability is itself uncertain, and completely
unknown — model it as $U \sim \mathrm{Uniform}(0,1)$, every value equally
plausible. Push it to the log-odds scale:

\[
  X \;=\; \operatorname{logit}(U).
\]

Uniform noise on $(0,1)$ goes in; some distribution over all of $\mathbb{R}$
comes out. Which one?

\begin{parts}
  \item Verify that $\operatorname{logit}$ and $\sig$ really are inverse:
        show $\sig(\operatorname{logit}(u)) = u$. \pt{1}
  \item Find the CDF of $X$ directly. Start from $F_X(x) = \PP(X \le x)$,
        rearrange $\log\frac{U}{1-U} \le x$ into a statement about $U$ alone, and
        then use $\PP(U \le c) = c$ for a uniform. \pts{3}
  \item Differentiate to get the density $f_X$. Show it can be written in all
        three of these ways, and check it integrates to $1$:
        \[
          f_X(x) \;=\; \sig(x)\big(1 - \sig(x)\big)
                  \;=\; \frac{e^{-x}}{(1+e^{-x})^{2}}
                  \;=\; \frac{1}{2\big(\cosh x + 1\big)} .
        \]
        \pts{3}
  \item Show $f_X(-x) = f_X(x)$, and deduce $\EE[X]$ without integrating. \pts{2}
  \item This distribution has a name — state it. Then explain in one or two
        sentences why the answer to (b) was a foregone conclusion before you did
        any algebra. \pt{1}
\end{parts}

\noindent\textit{Part (e) is the real content. Feeding a uniform through the
inverse of a CDF is precisely how one samples from that CDF — so applying
$\sig^{-1}$ to a uniform could only ever produce the law whose CDF is $\sig$.
This is the same ``sampling is transport'' idea the course uses when it writes a
sample as $Q(U)$.}

\total{10}

\bigskip
\hrule
\bigskip

\section*{Problem 3}

A website is testing a new checkout page. It is shown to $n$ visitors, chosen
independently of one another, and each either completes a purchase or does not.
The conversion probability $p \in (0,1)$ is unknown — estimating it is the entire
point of running the test.

Record $x_i = 1$ if visitor $i$ converts and $0$ otherwise, so
$x_1,\dots,x_n$ are i.i.d.\ $\mathrm{Bernoulli}(p)$, and let

\[
  S \;=\; \sum_{i=1}^{n} x_i
\]

be the total number of conversions. Notice that the raw data is the whole
sequence of $0$s and $1$s, but the number you would actually report is $S$.

\begin{parts}
  \item Write $\PP(x_i = x)$ for $x \in \{0,1\}$ as a \emph{single} expression in
        $p$ and $x$, with no cases. (This one line makes part (c) painless.)
        \pt{1}
  \item Show that $S \sim \mathrm{Binomial}(n,p)$, that is
        $\PP(S = k) = \binom{n}{k}p^{k}(1-p)^{n-k}$. Your argument needs three
        steps: why any one specific sequence with $k$ conversions has probability
        $p^{k}(1-p)^{n-k}$; how many such sequences there are; and why you are
        allowed to add their probabilities. \pts{3}
  \item Write the likelihood $L(p)$ of the observed sequence and the
        log-likelihood $\ell(p)$. \pts{2}
  \item Derive the maximum likelihood estimator $\hat p$ from $\ell'(p) = 0$, and
        confirm it is a maximum using $\ell''(p)$. Is the answer the one your
        intuition would have guessed? \pts{3}
  \item Does anything change if you start instead from the Binomial likelihood
        for $S$, rather than the Bernoulli likelihood for the whole sequence?
        Say what happens to $\binom{n}{k}$, and why. \pt{1}
\end{parts}

\total{10}

\bigskip
\hrule
\bigskip

\section*{Problem 4}

A lab buys a second-hand instrument. Its datasheet is missing, so nobody knows
the range over which it reports: readings are known to be spread evenly across
some interval, but the endpoints of that interval are part of what has been lost.

Model $n$ readings as i.i.d.\ from $\mathrm{Uniform}(a,b)$ with $a < b$ both
unknown, and recover the calibration from the data.

\begin{parts}
  \item Write down the likelihood $L(a,b)$ for the whole sample. Take care: the
        density equals $1/(b-a)$ only \emph{inside} the interval and $0$ outside,
        so your expression must record that every observation lies in $[a,b]$.
        \pts{3}
  \item Explain why setting $\partial \ell/\partial a = 0$ and
        $\partial \ell/\partial b = 0$ does \emph{not} work here. \pts{2}
  \item Argue directly from the shape of $L$ — not by differentiating — that
        \[
          \hat a = \min_i x_i, \qquad \hat b = \max_i x_i .
        \]
        \pts{3}
  \item The estimated interval $[\hat a, \hat b]$ is always contained inside the
        true $[a,b]$. What does that say about whether $\hat b - \hat a$ is an
        unbiased estimate of the true width $b-a$, and in which direction it goes
        wrong? Would collecting more readings help? \pts{2}
\end{parts}

\noindent\textit{Set this beside Problem 3. There, the maximum was found by
calculus. Here calculus is the wrong instrument entirely and the answer falls out
of the support constraint. Recognising which situation you are in is the skill
being tested.}

\total{10}

\bigskip
\hrule
\bigskip

\section*{Problem 5}

A postal sorting machine reads the handwritten digit on an envelope and must
decide which of the $K = 10$ classes $\{0,1,\dots,9\}$ it is. From the image
$x_i$ a network produces not one number but a \emph{vector} of $K$ scores,

\[
  z_i \;=\; \big(z_{i1}, \dots, z_{iK}\big) \;=\; f_\theta(x_i) \;\in\; \mathbb{R}^{K},
\]

one score per class, unconstrained in sign or size. These are called
\textbf{logits}. To report probabilities they are passed through the
\textbf{softmax}

\[
  q_{ik} \;=\; \frac{e^{z_{ik}}}{\sum_{m=1}^{K} e^{z_{im}}},
  \qquad k = 1,\dots,K .
\]

The true label is written \textbf{one-hot}: $y_i = (y_{i1},\dots,y_{iK})$ with
$y_{ik} = 1$ for the correct class and $0$ for the other $K-1$. So exactly one
entry is a $1$, and $\sum_k y_{ik} = 1$.

The machine is trained by minimising the \textbf{cross-entropy loss}

\[
  \mathrm{CE}(\theta) \;=\; -\frac{1}{n}\sum_{i=1}^{n}\sum_{k=1}^{K}
    y_{ik}\,\log q_{ik} .
\]

Like every loss in this course, it is not a formula to be accepted. It is a
maximum-likelihood estimator wearing a different hat, and this problem takes the
hat off.

\begin{parts}
  \item \textbf{Softmax is a distribution.} Show that $q_{ik} > 0$ for every $k$
        and that $\sum_{k=1}^{K} q_{ik} = 1$, so the output really is a pmf over
        the $K$ classes. Then show that softmax is \emph{shift-invariant}: adding
        the same constant $c$ to every logit leaves $q_i$ unchanged. What does
        that say about how many of the $K$ logits are genuinely free? \pts{2}
  \item \textbf{The sampling distribution.} State the distribution of the label
        $y_i$ given the image $x_i$ that the machine is implicitly assuming, with
        $q_i$ as its parameter vector. Write its pmf as a \emph{single}
        expression in $y_i$ and $q_i$, with no cases and no summation over
        classes --- the one-hot encoding is what makes this possible. \pts{2}
  \item \textbf{The generative model, and its assumptions.} Write the likelihood
        of the whole labelled dataset under that assumption. Take the logarithm,
        negate, divide by $n$, and show you land exactly on $\mathrm{CE}(\theta)$.
        Then state the three modelling assumptions you used: one about the
        \emph{form} of each individual label's distribution, one about the
        relationship \emph{across} the $n$ envelopes, and one about what role the
        images $x_i$ themselves play. \pts{3}
  \item \textbf{The gradient.} For a single envelope, write $\ell(z) = -\sum_k
        y_k \log q_k$ with $q = \mathrm{softmax}(z)$. Show first that
        \[
          \frac{\partial \log q_k}{\partial z_j} \;=\; \delta_{kj} - q_j,
          \qquad
          \delta_{kj} = \begin{cases} 1 & k = j\\ 0 & k \ne j\end{cases}
        \]
        and deduce
        \[
          \frac{\partial \ell}{\partial z_j} \;=\; q_j - y_j ,
        \]
        that is, predicted minus observed --- in every coordinate at once. Say
        where in your derivation you used $\sum_k y_k = 1$. \pts{2}
  \item \textbf{Back to two classes.} Set $K = 2$. Show that
        $q_2 = \sig(z_2 - z_1)$, and hence that the loss above collapses to the
        binary cross-entropy of the lecture notes with a single logit
        $z = z_2 - z_1$. Which part of (a) explains why two logits were one too
        many? \pt{1}
\end{parts}

\noindent\textit{Part (b) deserves the most thought. The machine is not
generating envelopes --- it is generating \emph{labels}, conditionally on images
it takes as given. That is why what you write down is a conditional sampling
distribution, and why the likelihood that follows is a conditional one. Part (d)
explains why this loss is so pleasant to optimise: the gradient is just the
error, which is exactly the form the denoiser objective takes later in the
course.}

\total{10}

\end{document}
